\section{زنجیره شبیه‌سازی و نمونه‌ها}
\label{sec:sim}
ابتدا فرایند سخت را با \lr{MadGraph 3.5.3} در سطح پارتون تولید می‌کنیم. در این مرحله برش‌های زیر را اعمال می‌کنیم. این برش‌ها را طوری گذاشته‌ایم که جت‌های بوست‌شده‌ای تولید شود که هر دو کوارک \bq{} در آن جای بگیرند. وقتی \Zbos{} با تکانه عرضی بزرگ تولید می‌شود، زاویه بین دو محصول واپاشی کوچک می‌شود: تقریباً $\Delta R\approx 2m_\mathrm{Z}/\pt\approx 0.45$ برای \lr{$\pt=400$~GeV}. این مقدار از شعاع جت ($R=0.8$) کوچک‌تر است و هر دو \bq{} در یک جت می‌افتند. به همین دلیل در \lr{zj} و \ttZ روی \pt بوزون \Zbos{} برش \lr{$\pt(\mathrm{Z})>400$~GeV} گذاشته‌ایم. علاوه بر این برش‌ها، برش $|\eta|<2$ روی جت‌های بازسازی‌شده اعمال می‌شود.

پس‌زمینه‌ها هم باید در همان ناحیه بوست‌شده تولید شوند، وگرنه مقایسه با سیگنال منصفانه نیست و بیشتر رویدادها بعد از برش‌های تحلیل از بین می‌روند. در \ttbb برش روی \pt زوج \bb (بزرگ‌تر از ۴۰۰~\lr{GeV}) همین نقش را دارد، و شرط \lr{$\Delta R(\mathrm{b\bar{b}})<0.8$} دو کوارک \bq{} را به هم نزدیک می‌کند تا هر دو در یک جت با $R=0.8$ جای بگیرند و هندسه‌ای شبیه واپاشی \Zbos{} داشته باشیم.

همچنین برش‌های تحلیل، یعنی \lr{$\pt>400$~GeV} و \lr{$|\eta|<2$}، روی جت‌های بازسازی‌شده اعمال می‌شوند، چون پاسخ آشکارساز و تابش بیرون از مخروط جت دقیقاً یکی نیستند؛ بخشی از جت‌ها بعد از برش \pt زیر ۴۰۰ می‌افتند (جدول~\ref{tab:cutflow}).

سپس با \lr{MadSpin} ذرات سنگین را واپاشی می‌دهیم. در \lr{zj} و \ttZ بوزون \Zbos{} به یک زوج \bb واپاشی می‌کند. در \ttZ هر یک از دو کوارک تاپ به یک کوارک \bq{} و یک بوزون \lr{W} می‌پاشد و هر \lr{W} به یک لپتون باردار و نوترینوی متناظر آن؛ لپتون باردار الکترون یا میون است (در رویدادهای بررسی‌شده تقریباً نیمی الکترون و نیمی میون‌اند و لپتون تاو دیده نمی‌شود). در \ttbb هم دو تاپ دقیقاً به همین شکل، یعنی دولپتونی، واپاشی داده می‌شوند. \lr{MadSpin} همبستگی‌های اسپینی بین محصولات این واپاشی‌ها را حفظ می‌کند.

بعد از آن \lr{Pythia 8.312} بارش پارتونی و هادرونیزاسیون را انجام می‌دهد. در بارش پارتونی، کوارک‌ها و گلوئون‌های اولیه با تابش گلوئون و شکافت گلوئون به زوج کوارک، به تعداد زیادی پارتون با انرژی کمتر تبدیل می‌شوند. هادرونیزاسیون (\lr{hadronization}) مرحله بعد است: به دلیل محصورشدگی رنگ، کوارک‌ها و گلوئون‌ها نمی‌توانند به‌صورت آزاد وجود داشته باشند؛ در پایان بارش، آنها به هادرون‌های بی‌رنگ (مزون‌ها و باریون‌ها) تبدیل می‌شوند، که در \lr{Pythia} با مدل رشته لوند~\cite{Andersson:1983ia} انجام می‌شود. هادرون‌های ناپایدار سپس واپاشی می‌کنند تا ذرات پایانی بماند.

\lr{Delphes 3.5.1pre09} پاسخ آشکارساز را با تنظیمات موجود در آزمایش \lr{CMS} شبیه‌سازی می‌کند و خروجی آن ردهای باردار و اثرات آشکارشده در کالریمترها (فوتون و هادرون خنثی) است. در پایان \lr{FastJet 3.4.1} این ذره‌ها را به جت تبدیل می‌کند: جت‌ها با الگوریتم \lr{anti-$k_\mathrm{T}$} و $R=0.8$ ساخته می‌شوند، که ذره‌های سخت را اول با ذره‌های نرم اطرافشان ادغام می‌کند و جت‌های گرد و مخروطی می‌دهد؛ توضیح مرحله‌به‌مرحله آن در پیوست~\ref{app:antikt} آمده است. در \lr{FastJet} فقط جت‌هایی با \lr{$\pt>20$~GeV} ذخیره می‌شوند و بعد متغیرهای زیرساختار و ویژگی‌های ذره‌ها محاسبه و ذخیره می‌شوند.

برای هر کار خروجی هر مرحله (\lr{MadGraph}، \lr{MadSpin}، \lr{Pythia} و \lr{Delphes}) ذخیره شده و یک بذر تصادفی (\lr{random seed}) یکسان برای همه این مراحل به‌کار می‌رود. بذر هر کار از شناسه نمونه و شماره کار ساخته می‌شود و در یک پرونده ثبت می‌شود، پس کل زنجیره تکرارپذیر است و با همان بذر دقیقاً همان رویدادها دوباره به دست می‌آیند (\lr{FastJet} تصادفی نیست و به بذر نیازی ندارد). همه نمونه‌ها در مرتبه درختی (\lr{LO}) تولید شده‌اند و هر دو نمونه پس‌زمینه با فرایند \lr{QCD} خالص: \bbbb با \lr{QED=0}، یعنی دامنه بدون ثابت جفت‌شدگی الکتروضعیف، و \ttbb با \lr{WEIGHTED<=4}. \lr{WEIGHTED} مرتبه وزن‌دار ثابت‌های جفت‌شدگی در دامنه است که در \lr{MadGraph} هر توان $g_s$ را ۱ و هر توان $e$ را ۲ حساب می‌کند. دامنه $\mathrm{gg}\to t\bar t b\bar b$ فقط $g_s^4$ دارد، یعنی وزن ۴، و هر دیاگرام با $e$ وزن بیشتر از ۴ می‌شود و کنار می‌رود.

هنگام بازسازی جت‌ها، با توجه به مادر هر ذره‌ای که داخل جت قرار دارد، برای هر جت یک سری برچسب ثبت می‌کنیم که نشان می‌دهد جت از چه ذره‌هایی ساخته شده است، یعنی منشأ جت. جتی را که هر دو کوارک \bq{} حاصل از \Zbos{} داخل آن باشد جت کامل می‌نامیم و جتی را که فقط یکی از آن دو داخل آن باشد جت جزئی. در نهایت، در هر چهار فرایند این جت‌ها را بر اساس برچسب‌هایشان به سیگنال و پس‌زمینه تقسیم می‌کنیم: جت‌های کامل (و در حالتی دیگر جت‌های کامل و جزئی با هم) در دو نمونه سیگنال به‌عنوان سیگنال برداشته می‌شوند و جت‌هایی که \bq{} غیرتاپی دارند در دو نمونه پس‌زمینه به‌عنوان پس‌زمینه.

خلاصه‌ای از برش‌ها و برچسب‌های منشأ را در جدول~\ref{tab:samples} و جدول~\ref{tab:origin} می‌توانید ببینید.
\begin{table}[!ht]
\centering
\caption{نمونه‌ها و برش‌های سطح پارتون. در هر نمونه ۲۰ کار با ۱۰۰۰ رویداد، مجموعاً ۲۰۰۰۰ رویداد، اجرا شده است.}
\label{tab:samples}
\begin{tabular}{llp{6.5cm}}
\toprule
نمونه & نقش & برش سطح پارتون \\
\midrule
\lr{zj} & سیگنال & \lr{$\pt(\mathrm{Z})>400$~GeV}، بدون برش $\eta$ \\
\ttZ & سیگنال & \lr{$\pt(\mathrm{Z})>400$~GeV}، بدون برش $\eta$ \\
\bbbb & پس‌زمینه & چهار \bq{} به دو زوج با کمترین مجموع $\Delta R$ تقسیم می‌شوند و برای هر دو زوج \lr{$\pt>400$~GeV} و \lr{$\Delta R<0.8$} \\
\ttbb & پس‌زمینه & \lr{$\pt(\mathrm{b\bar{b}})>400$~GeV} و \lr{$\Delta R(\mathrm{b\bar{b}})<0.8$} \\
\bottomrule
\end{tabular}
\end{table}

\begin{table}[!ht]
\centering
\caption{برچسب‌های منشأ جت.}
\label{tab:origin}
\begin{tabular}{lp{10cm}}
\toprule
برچسب & معنی \\
\midrule
\lr{z\_both} & هر دو \bq{} حاصل از \Zbos{} داخل جت‌اند (جت کامل) \\
\lr{z\_part} & فقط یکی از دو \bq{} حاصل از \Zbos{} داخل جت است (جت جزئی) \\
\lr{nb\_notop} & تعداد \bq{}های داخل جت که از تاپ نیامده‌اند \\
\lr{nb\_top} & تعداد \bq{}های داخل جت که از تاپ آمده‌اند \\
\lr{n\_hard\_g}, \lr{n\_hard\_q} & تعداد گلوئون و کوارک سخت داخل جت \\
\lr{d\_z}, \lr{d\_top} & فاصله جت از \Zbos{} و از تاپ \\
\bottomrule
\end{tabular}
\end{table}
